% Figure 4: Markov closure centerpiece (Session 18, Experiment B1).
% Source figure: outputs/session18/figures/figure4_markov_closure_centerpiece.pdf
% Source data:   outputs/session18/exp_b1_test3/physical_closure_noBN_unified.csv (864 rows)
% Aggregation:   z_markov mode, horizon H=16 frames after impact, mean absolute error,
%                bootstrap 95\,\% CI over encounters (n=42 Test B, n=24 Test C).

\begin{figure}
  \centering
  \includegraphics[width=\textwidth]{outputs/session18/figures/figure4_markov_closure_centerpiece.pdf}
  \caption[Markov closure of physical observables across eight latent baselines.]{%
    Markov-only forecast of physical observables at horizon $H=16$ frames after
    vortex impact, comparing JEPA against the Fukami AE and POD baseline
    families under a single unified transformer predictor (no output BatchNorm)
    fit on \texttt{train}. The panel grid is $2\times3$: rows distinguish
    \emph{Test B} (in-distribution, $n=42$ encounters, top) from \emph{Test C}
    ($G=+4$ out-of-distribution, $n=24$, bottom); columns report the lift
    coefficient $C_L$, the wake $y$-impulse $I_y$, and the wake enstrophy
    $\Omega_w$. Within each panel the seven bars compare, left to right, Fukami
    AE at $d\in\{3,32,64\}$ (red shades), POD at $d\in\{16,32,64\}$ (blue
    shades), and JEPA $d=64$ (green, hatched); error bars are bootstrap 95\,\%
    CI ($n=2000$ resamples over encounters); horizontal dotted lines mark the
    DNS-oracle floor (best closure from ground-truth observables, also fit on
    \texttt{train}). JEPA wins decisively on the wake structure on both
    splits: on Test B $\Omega_w$ absolute error drops from
    $74.2$--$89.3$ (POD family) and $55.7$--$73.3$ (Fukami AE family) to $41.3$
    for JEPA $d=64$; on Test C it drops from $217.0$--$223.4$ (POD) and
    $233.6$--$262.6$ (Fukami AE) to $139.2$, the only entry below the
    DNS-oracle floor on this OOD split, and the same advantage carries over to
    circulation$_{\pm}$ in Table~\ref{tab:b1-headline} where JEPA reduces error
    by $1.5$--$3.3\times$ over POD. On the aerodynamic coefficients $C_L$ and
    $C_D$ JEPA ties the strongest classical baseline at matched capacity
    (matched-$d=64$ Test B $C_L$: $0.80$ for JEPA, $0.78$ for POD $d=64$ and
    Fukami AE $d=64$, with overlapping CIs); on $I_y$ POD wins on Test C, an
    honest task-dependent trade-off. Full headline numbers across
    \texttt{train}, \texttt{val}, \texttt{test\_b}, and \texttt{test\_c} with
    three-seed encoder variance, five-fold probe cross-validation, the six
    physical observables, and the horizon sweep $H\in\{8,16,32\}$ are reported
    in Table~\ref{tab:b1-headline}. Source data:
    \texttt{outputs/session18/exp\_b1\_test3/physical\_closure\_noBN\_unified.csv}
    (ridge probes; KRR-RBF sensitivity in Figure~S6).
  }
  \label{fig:b1-markov-closure}
\end{figure}
